% Complete Chapter 2 introductory paragraph, original PDF 21 / printed 11.
% block: homological-algebra-heading page:21 type:chapter
\clearpage\setcounter{section}{1}
\section{{Basic notions of homological algebra}}\label{chap:homological-algebra}

% block: homological-algebra-history page:21 type:ordinary-paragraph
{The basic structures of homological algebra, such as chain complexes and homology groups, were introduced in the 19th century with the aim of defining topological invariants such as the Betti numbers via the counting of discrete structures. However, their antecedents go back much further, to Euler and even Descartes. With the development of differential forms by E. Cartan at the start of the 20th century and the study of their cohomology by de Rham in the 1930s these same algebraic structures appeared in the context of spaces of functions acted on by partial differential operators. Later yet they became objects of study from an abstract algebraic point of view. Thus homological algebra has long played a fundamental role in algebraic topology, differential geometry, and algebra. As useful as it is, however, it is often unfamiliar to applied mathematicians and numerical analysts. Therefore, in this chapter, we include a short self-contained introduction to the basic aspects of homological algebra we shall need. The key points are summarized in Box~\ref{box:homological-algebra-summary} at the end of the chapter.}
